\documentclass[a4paper,12pt]{article}
\addtolength{\oddsidemargin}{-1.cm}
\addtolength{\textwidth}{2cm}
\addtolength{\topmargin}{-2cm}
\addtolength{\textheight}{3.5cm}
\newcommand{\HRule}{\rule{\linewidth}{0.5mm}}
\makeindex

\usepackage{longtable}
\usepackage[pdftex]{graphicx}
\usepackage{makeidx}
\usepackage{hyperref}
\usepackage{verbatim}
\hypersetup{
    colorlinks=true,
    linkcolor=blue,
    filecolor=magenta,      
    urlcolor=cyan,
}


% define the title
\author{IMAPKD}
\title{ Project Tender}
\begin{document}
\setlength{\parskip}{6pt}

% generates the title
\begin{titlepage}

\begin{center}
% Upper part of the page       
\includegraphics[width=1\textwidth]{./University_of_Pretoria_Logo.PNG}\\[0.4cm]  
\textsc{\LARGE COS 301}\\[0.9cm]
\textsc{\LARGE Department Of Computer Science}\\[0.3cm]

%Title

\HRule \\[0.4cm]
{ \huge \bfseries Architectural Requirements and  Initial Architecture Design Functional Requirements}\\[0.1cm]
\HRule \\[0.4cm]  

% Group Members 

\begin{minipage}{0.4\textwidth}
\begin{flushleft} \large

\emph{\Large Group Members:}\\[0.4cm]    
\emph{}\\
{\Large Diana {Obo}} \\
\emph{}\\
\emph{}\\
{\Large Priscilla {Madigoe}}\\
\emph{}\\
\emph{}\\
{\Large Kudzai {Muranga}} \\
\emph{}\\
\emph{}\\
{\Large Sandile {Khumalo}}\\
\emph{}\\
\emph{}\\

\end{flushleft}
\end{minipage}
\begin{minipage}{0.4\textwidth}
\begin{flushright} \large

\emph{ \Large Student numbers:} \\[0.4cm]  
\emph{}\\
{\Large u13134885}\\
\emph{}\\
\emph{}\\
{\Large u13049128}\\
\emph{}\\
\emph{}\\
{\Large u13278012}\\
\emph{}\\
\emph{}\\
{\Large u12031748}\\
\emph{}\\
\emph{}\\

\end{flushright}
\end{minipage}

% bottom section of title page

{\large \today}
\end{center}
\end{titlepage}
\renewcommand{\thesection}{\arabic{section}}

\newpage
\begin{center}
\textsc{\Large IMPAKD link}\\[0.5cm]
For further references see \href{https://github.com/u13278012/IMPAKD/}{gitHub}.
\today
\end{center}
\newpage
\tableofcontents{}

\newpage
\section{Vision}

\newpage
\section{Background}

\newpage
\section{Software Architecture}
\subsection{Architecture requirements}
\subsubsection{Architectural scope}
\subsubsection{Quality requirements}
%Priscilla's Section - Start

\subsubsection{Integration and access channel requirements}
	\begin{itemize}
		\item{\bfseries \underline{Integration}}
		\begin{itemize}
           		\item Logging into the system is done over a HTTPS POST method.
			\item The user's login details are kept in an HTTP session so the user does not need to log in everytime he/she 					makes a request to the server.
			\item The HTTPS sessions are invalidated when the user terminates his/her session by logging out.
			\item Communication between the server (back-end) and the webpage (front-end) will be facilitated by the REST 					method which uses JSON objects and HTTPS methods to send requests and get responses.
			\item The "Create, Read, Update and Delete" or CRUD actions that will make changes to the databse will be 					logged automatically. This will ensure auditability of the system. 
		\end{itemize}

		\item{\bfseries \underline{Human Access Channel}}
			\begin{itemize}
				\item End-users interact with the Web client to display the required information and do desired actions. 
			\end{itemize}

		\item{\bfseries \underline{System Access Channel}}
			\begin{itemize}
				\item The Web-based component of the system will be implememted in "Ember.js" which utilises JavaScript, 				HTML and "Handlebar.js".  
			\end{itemize}
	\end{itemize}
%Priscilla's Section - End

\subsubsection{Architectural constraints}

% Architectural patterns Section - start
\subsection{Architectural patterns or styles}
\textbf{MVC (\textit{Model View Controller})}\\
Allows the system's states to change and it encapsulates the interactions
from the user and transforms these intercations into business logic.\\
\textsc{Reason:}
	\begin{itemize}
		\item Reduce presentation layers complexity and  improves flexibility
			\begin{itemize}
				\item Separates responsibilities
					\begin{itemize}
						\item Provide view onto information -\textbf{\textit{ View}}
						\item React to user events - \textbf{\textit{Controller}}
						\item Provide business services  and data - \textbf{\textit{Model}}
					\end{itemize}
				\item Allows each component to change independently
			\end{itemize}
		\item Full decoupling
			\begin{itemize}
				\item Model from both \textit{view} and \textit{controller}
			\end{itemize}
		\item Simplification
			\begin{itemize}
				\item Through separation of concerns
			\end{itemize}
		\item Reuse
			\begin{itemize}
				\item \textit{Model} components and \textit{View} components 
			\end{itemize}
		\item Maintainability
			\begin{itemize}
				\item Different components can be used, developed and maintained by different memebers of a team 
					\begin{itemize}
						\item \textit{Model} - backened developers
						\item \textit{View} - UI designers
						\item \textit{Controller} - Front-end developers
					\end{itemize}
			\end{itemize}
		\item Improved Testability
			\begin{itemize}
				\item Model/business services tested independently of UI
				\item UI tested with mock model
			\end{itemize}
	\end{itemize}

%Architectural patterns Section - end

\subsection{Architectural tactics or strategies}
\subsection{Use of reference architectures and frameworks}
\subsection{Access and integration channels}
\subsection{Technologies}

\newpage
\section{Functional requirements and application design}

% Use case prioritization Section - start
\subsection{Use case prioritization}
 \underline{\textbf{Critical:}}
	\begin{itemize}
		\item calculateROI
		\item getDefaultValues
		\item setDeafultValues
	\end{itemize}
\underline{\textbf{important:}}
	\begin{itemize}
		\item Register
		\item Login
		\item logout
		\item addProperty
		\item updateProperty	
		\item deleteProperty
		\item displayGraphs
		\item displayStatistics	
	\end{itemize}
\underline{\textbf{Nice-to-have:}}
	\begin{itemize}
		\item updateProfile
		\item generateReport	
	\end{itemize}

% Use case prioritization Section - start

\subsection{Use case/Services contracts}
\subsection{Required functionality}
\subsection{Process specifications}
\subsection{Domain Model}

\newpage
\section{Open Issues}



\end{document}

